% 图 2-7 三种实现方案的调用路径与周期（坐标单位 cm）
\begin{tikzpicture}[font=\footnotesize,
  b/.style={draw=softline, rounded corners=1.5pt, fill=macblue!55, align=center, inner sep=1.8pt, font=\scriptsize, minimum height=0.85cm},
  w/.style={draw=softline, rounded corners=1.5pt, fill=white, align=center, inner sep=1.8pt, font=\scriptsize, minimum height=0.85cm},
  lab/.style={font=\scriptsize, align=left, text width=2.2cm},
  ar/.style={-{Stealth[length=1.4mm]}, thick}]

\node[font=\scriptsize\bfseries, anchor=west] at (-7.7,1.15) {方案};
\node[font=\scriptsize\bfseries] at (-0.6,1.15) {调用路径};
\node[font=\scriptsize\bfseries, anchor=west] at (5.9,1.15) {当前加密周期};

% 一、纯软件
\node[lab, anchor=west] at (-7.7,0.3) {一、纯软件\\（SW）};
\node[b, minimum width=2.2cm, fill=macblue!70] (a1) at (-3.9,0.3) {RV32\\指令流};
\node[w, minimum width=2.4cm, fill=macblue!25] (a2) at (-1.0,0.3) {sbox / 线性层\\逐轮软件计算};
\node[w, minimum width=1.9cm, fill=macblue!25] (a3) at (2.0,0.3) {栈与\\BRAM};
\node[b, minimum width=2.0cm, fill=macblue!70] (a4) at (4.6,0.3) {全部由\\CPU 执行};
\node[font=\scriptsize, anchor=west] at (5.9,0.3) {\textbf{23\,342}};
\foreach \x/\y in {a1/a2,a2/a3,a3/a4} {\draw[ar] (\x) -- (\y);}

% 二、MMIO
\node[lab, anchor=west] at (-7.7,-1.3) {二、内存映射\\（MMIO）};
\node[b, minimum width=2.2cm, fill=macgreen] (b1) at (-3.9,-1.3) {CPU 写\\寄存器};
\node[w, minimum width=2.4cm, fill=macgreen!40] (b2) at (-1.0,-1.3) {SoC MMIO\\寄存器事务};
\node[w, minimum width=1.9cm, fill=macgreen!40] (b3) at (2.0,-1.3) {\texttt{ascon\_mmio}\\置换核};
\node[b, minimum width=2.0cm, fill=macgreen] (b4) at (4.6,-1.3) {CPU 轮询\\done};
\node[font=\scriptsize, anchor=west] at (5.9,-1.3) {\textbf{1\,439}};
\foreach \x/\y in {b1/b2,b2/b3,b3/b4} {\draw[ar] (\x) -- (\y);}

% 三、ISE
\node[lab, anchor=west] at (-7.7,-2.9) {三、定制指令\\（ISE）};
\node[b, minimum width=2.2cm, fill=macamber] (c1) at (-3.9,-2.9) {\texttt{ascon.p12}\\单条指令};
\node[w, minimum width=2.4cm, fill=macamber!45] (c2) at (-1.0,-2.9) {EX 级直连\\不经总线};
\node[w, minimum width=1.9cm, fill=macamber!45] (c3) at (2.0,-2.9) {\texttt{ascon\_ise}\\320-bit 私有态};
\node[b, minimum width=2.0cm, fill=macamber] (c4) at (4.6,-2.9) {hold / ack\\流水推进};
\node[font=\scriptsize, anchor=west] at (5.9,-2.9) {\textbf{3\,995}};
\foreach \x/\y in {c1/c2,c2/c3,c3/c4} {\draw[ar] (\x) -- (\y);}

\node[draw=softline, rounded corners=2pt, fill=macgray, align=left, font=\scriptsize, text width=14.6cm, anchor=north west]
  at (-7.7,-3.7) {当前 PYNQ release，KAT \#529，PT16 B/AD0，各预热 5 次后 30 次采样：
  MMIO $16.221\times$、ISE $5.843\times$。
  周期为 RV32 固件 AEAD 调用边界，不含 PS/UART；PS 调用墙钟必须单列。};
\end{tikzpicture}
